\section{Comments}

All functions must include Doxygen comments. This section establishes comprehensive commenting standards for the STAR firmware codebase.

\subsection{General Commenting Guidelines}

\begin{itemize}
    \item \textbf{Use Block Comments (\texttt{/* */})}: Always prefer block comments over line comments (\texttt{//}).
    \item \textbf{Explain Why, Not What}: Comments should explain the reasoning behind code, not what the code does (which should be clear from the code itself).
    \item \textbf{Keep Comments Updated}: Outdated comments are worse than no comments.
    \item \textbf{English Only}: All comments must be in English.
\end{itemize}

\begin{lstlisting}[language=C]
/* CORRECT - Block comment style */
/* Initialize the mutex for thread-safe access */
manager->mutex = xSemaphoreCreateMutex();

/* WRONG - Line comment style */
// Initialize the mutex for thread-safe access
manager->mutex = xSemaphoreCreateMutex();
\end{lstlisting}

\subsection{File Header Comments}

Every source and header file should begin with a file path comment, followed by include guards (for headers) and a \texttt{@file} Doxygen block for headers with significant content.

\begin{lstlisting}[language=C]
/* lib/star_bus/include/star_bus_manager.h */

#ifndef STAR_COMPONENT_BUS_MANAGER_H
#define STAR_COMPONENT_BUS_MANAGER_H

#ifdef __cplusplus
extern "C" {
#endif

#include "esp_err.h"
#include "star_bus_manager_types.h"

/**
 * @file star_bus_manager.h
 * @brief Unified bus manager for I2C, SPI, GPIO, and other bus protocols
 *
 * The bus manager provides a central registry for managing multiple bus
 * configurations. It supports dependency injection of error handlers and
 * pin validators for loose coupling.
 *
 * Key Features:
 * - Thread-safe bus management with mutex protection
 * - Support for multiple bus types (I2C, SPI, GPIO, DHT22, etc.)
 * - Automatic pin registration and conflict detection
 * - Safe bus access via callback pattern (prevents use-after-free)
 *
 * Example Usage:
 * @code
 * // === Basic Bus Manager Setup ===
 *
 * #include "star_bus_manager.h"
 * #include "star_bus_config.h"
 * #include "star_error_handler.h"
 *
 * // 1. Create error handler and pin validator interfaces
 * error_handler_t error_handler;
 * error_handler_init(&error_handler, 3, 100, 5000, NULL, NULL);
 *
 * star_error_interface_t error_iface;
 * star_pin_interface_t pin_iface;
 * error_handler_get_interface(&error_iface, &error_handler);
 * pin_validator_get_interface(&pin_iface);
 *
 * // 2. Initialize bus manager with dependency injection
 * star_bus_manager_t bus_manager;
 * star_bus_manager_init(&bus_manager, "main", &error_iface, &pin_iface);
 * @endcode
 */
\end{lstlisting}

\subsection{Source File Header}

Source files begin with a file path comment.

\begin{lstlisting}[language=C]
/* lib/star_bus/src/star_bus_manager.c */

#include "star_bus_manager.h"

#include "freertos/FreeRTOS.h"
#include "freertos/semphr.h"
/* ... */
\end{lstlisting}

\subsection{Doxygen Documentation}

\subsubsection{Function Documentation}

All public functions must have Doxygen documentation in the header file.

\begin{lstlisting}[language=C]
/**
 * @brief Initialize a bus manager instance.
 *
 * Allocates resources (like a mutex) for the manager. Must be called
 * before any other manager functions.
 *
 * @param[out] manager     Pointer to bus manager structure to initialize.
 *                         Must not be NULL.
 * @param[in]  tag         Optional tag string for logging associated with
 *                         this manager instance. If NULL or empty, a
 *                         default tag is used.
 * @param[in]  error_iface Optional error handler interface (can be NULL).
 * @param[in]  pin_iface   Optional pin validator interface (can be NULL).
 * @return esp_err_t ESP_OK on success, ESP_ERR_INVALID_ARG if manager
 *                   is NULL, ESP_ERR_NO_MEM if mutex creation fails.
 */
esp_err_t star_bus_manager_init(star_bus_manager_t*     manager,
                                const char*             tag,
                                star_error_interface_t* error_iface,
                                star_pin_interface_t*   pin_iface);
\end{lstlisting}

\subsubsection{Parameter Direction Markers}

Use \texttt{[in]}, \texttt{[out]}, or \texttt{[in,out]} to indicate parameter direction.

\begin{lstlisting}[language=C]
/**
 * @brief Write data to an I2C device register.
 *
 * @param[in]  config        Pointer to I2C bus configuration.
 * @param[in]  data          Pointer to data buffer to write.
 * @param[in]  len           Number of bytes to write.
 * @param[in]  reg_addr      Target register address.
 * @param[out] bytes_written Pointer to store actual bytes written.
 *                           Can be NULL if not needed.
 * @return esp_err_t ESP_OK on success.
 */
static esp_err_t internal_star_bus_i2c_write(const star_bus_config_t* config,
                                              const uint8_t*           data,
                                              size_t                   len,
                                              uint8_t                  reg_addr,
                                              size_t*                  bytes_written);

/**
 * @brief Record an error and manage retry/reset logic
 *
 * @param[in,out] handler     Pointer to error handler structure
 * @param[in]     error       Error code that occurred
 * @param[in]     description Human-readable description of the error
 * @param[in]     file        Source file where error occurred
 * @param[in]     line        Line number where error occurred
 * @param[in]     func        Function name where error occurred
 * @return ESP_OK if error was handled, or the original error code.
 */
esp_err_t error_handler_record_error(error_handler_t* handler,
                                     esp_err_t        error,
                                     const char*      description,
                                     const char*      file,
                                     int32_t          line,
                                     const char*      func);
\end{lstlisting}

\subsubsection{Code Examples with @code/@endcode}

Include usage examples in documentation blocks.

\begin{lstlisting}[language=C]
/**
 * @file star_error_interface.h
 * @brief Error handler interface for dependency inversion
 *
 * Example Usage:
 * @code
 * // === Creating and Using Error Interface ===
 *
 * #include "star_error_interface.h"
 * #include "star_error_handler.h"
 *
 * // 1. Create concrete error handler
 * error_handler_t error_handler;
 * error_handler_init(&error_handler, 3, 100, 5000, NULL, NULL);
 *
 * // 2. Get interface from implementation
 * star_error_interface_t error_iface;
 * error_handler_get_interface(&error_iface, &error_handler);
 *
 * // 3. Inject into components that need error handling
 * star_bus_manager_init(&bus_manager, "main", &error_iface, &pin_iface);
 *
 * // 4. Use the interface
 * if (error_iface.can_retry(error_iface.ctx)) {
 *     // Retry operation...
 * }
 * @endcode
 */
\end{lstlisting}

\subsubsection{Struct Member Documentation}

Document struct members inline with Doxygen.

\begin{lstlisting}[language=C]
/**
 * @brief Bus configuration structure (common part).
 */
typedef struct star_bus_config {
    const char*     name;        /**< Unique name for this bus/device instance */
    star_bus_type_t type;        /**< Type of bus (I2C, SPI, etc.) */
    bool            initialized; /**< Whether this bus has been initialized */
    void*           handle;      /**< Driver/device handle */
    void*           user_ctx;    /**< User context pointer for callbacks */

    union {
        star_i2c_bus_config_t  i2c;  /**< I2C-specific configuration */
        star_spi_bus_config_t  spi;  /**< SPI-specific configuration */
        star_gpio_bus_config_t gpio; /**< GPIO-specific configuration */
        star_adc_bus_config_t  adc;  /**< ADC-specific configuration */
    } proto;

    struct star_bus_config* next; /**< Next bus config in linked list */
} star_bus_config_t;
\end{lstlisting}

\subsection{Special Comment Keywords}

Use these keywords to mark special comments for tracking.

\begin{lstlisting}[language=C]
/* TODO: Implement retry logic for I2C bus reset */

/* FIXME: Memory leak when bus removal fails mid-operation */

/* BUG: Race condition possible if mutex timeout is too short */

/* XXX: This code assumes little-endian byte order */

/* NOTE: ESP-IDF uses mosi_io_num internally - we map to COPI here */

/* WARNING: Must hold mutex before calling this function */
\end{lstlisting}

\subsubsection{Keyword Definitions}

\begin{itemize}
    \item \textbf{TODO}: Work that needs to be done but is not urgent.
    \item \textbf{FIXME}: Known bugs or issues that need to be fixed.
    \item \textbf{BUG}: Confirmed bugs in the code.
    \item \textbf{XXX}: Dangerous or hacky code that works but should be reviewed.
    \item \textbf{NOTE}: Important information about the code.
    \item \textbf{WARNING}: Critical information that must be understood before modifying.
\end{itemize}

\subsection{Section Comments}

Use section comments to organize code within files.

\begin{lstlisting}[language=C]
/* === Includes === */
#include "star_bus_manager.h"
#include "freertos/FreeRTOS.h"

/* === Constants === */
static const char* s_TAG = "STAR_BUS_MANAGER";

/* === Private Function Prototypes === */
static esp_err_t internal_register_bus_pin(...);

/* === Public Functions === */
esp_err_t star_bus_manager_init(...) { /* ... */ }

/* === Private Functions === */
static esp_err_t internal_register_bus_pin(...) { /* ... */ }

/* --- Pin Registration Helpers --- */
/* Subsection within a section */
\end{lstlisting}

\subsection{Aligned Comments}

Align comments after declarations for improved readability.

\begin{lstlisting}[language=C]
typedef struct error_handler {
    uint32_t          error_count;         /**< Total errors recorded */
    uint32_t          max_retries;         /**< Max retry attempts */
    uint32_t          current_retry;       /**< Current retry counter */
    uint32_t          base_retry_delay;    /**< Base delay (ms) */
    uint32_t          max_retry_delay;     /**< Max delay (ms) */
    uint32_t          current_retry_delay; /**< Current delay with backoff */
    esp_err_t         last_error;          /**< Last recorded error code */
    bool              in_error_state;      /**< Whether in error state */
    void*             reset_context;       /**< Context for reset function */
    SemaphoreHandle_t mutex;               /**< Mutex for thread-safety */

    esp_err_t (*reset_fn)(void* context);  /**< Optional reset function */
} error_handler_t;
\end{lstlisting}

\subsection{Implementation Comments}

Comments within function bodies explain non-obvious logic.

\begin{lstlisting}[language=C]
esp_err_t star_bus_manager_add_bus(star_bus_manager_t* manager,
                                    star_bus_config_t*  config)
{
    /* Validate inputs using ESP-IDF macro */
    ESP_RETURN_ON_FALSE(manager && config, ESP_ERR_INVALID_ARG,
                        s_TAG, "Manager or config pointer is NULL");

    esp_err_t ret = ESP_OK;

    /* Acquire mutex with timeout to prevent deadlock */
    if (xSemaphoreTake(manager->mutex,
            pdMS_TO_TICKS(CONFIG_STAR_KCONFIG_BUS_MUTEX_TIMEOUT_MS))
            != pdTRUE) {
        ESP_LOGE(manager->tag, "Timeout acquiring mutex");
        return ESP_ERR_TIMEOUT;
    }

    /* Check for duplicate name - bus names must be unique */
    for (star_bus_config_t* curr = manager->buses; curr; curr = curr->next) {
        if (curr->name && strcmp(curr->name, config->name) == 0) {
            ret = ESP_ERR_INVALID_STATE;
            goto add_give_mutex;  /* Single exit point pattern */
        }
    }

    /* Initialize hardware before adding to list
     * This ensures manager's SPI host tracking is updated correctly */
    if (!config->initialized) {
        ret = star_bus_config_init(config, manager);
        if (ret != ESP_OK) {
            goto add_give_mutex;
        }
    }

    /* Add to head of list (O(1) insertion) */
    config->next = manager->buses;
    manager->buses = config;

add_give_mutex:
    xSemaphoreGive(manager->mutex);
    return ret;
}
\end{lstlisting}

\subsection{Private Function Documentation}

Static functions should have brief documentation in the source file.

\begin{lstlisting}[language=C]
/**
 * @brief Helper to register a single pin via the interface
 */
static esp_err_t internal_register_bus_pin(star_pin_interface_t* pin_iface,
                                            gpio_num_t            pin,
                                            const char*           bus_name,
                                            const char*           usage,
                                            bool                  can_be_shared)
{
    if (pin < 0 || pin >= GPIO_NUM_MAX) {
        return ESP_OK;  /* Not an error, pin is not used */
    }

    if (!pin_iface || !pin_iface->register_pin) {
        return ESP_OK;  /* No interface, skip registration */
    }

    /* Create description: "BusName: Usage" (e.g., "I2C0: SDA") */
    char desc[64];
    snprintf(desc, sizeof(desc), "%s: %s", bus_name, usage);

    return pin_iface->register_pin(pin_iface->ctx, pin, desc, can_be_shared);
}
\end{lstlisting}

\subsection{Summary}

\begin{itemize}
    \item Use block comments (\texttt{/* */}), never line comments (\texttt{//})
    \item Document all public functions with Doxygen in headers
    \item Use \texttt{[in]}, \texttt{[out]}, \texttt{[in,out]} for parameter direction
    \item Include \texttt{@code}/@\texttt{endcode} examples in file and complex function docs
    \item Use TODO, FIXME, BUG, XXX, NOTE, WARNING keywords for tracking
    \item Align comments after struct members for readability
    \item Use section comments to organize source files
\end{itemize}

